\begin{proof}[Proof of Theorem~\ref{main:birkhoff_unique}]
    Identify the vertices of $P(D)$ and $B_n$
    with $G$ and $S_n$, respectively.
    Let 
    $\gamma\colon G\to S_n$ be a combinatorial isomorphism.
    Then $\gamma$ induces an isomorphism $\kappa_{\gamma}$
    from the combinatorial symmetry group
    $A$ of $P(D)$ onto the combinatorial symmetry group
    $S_n\wr C_2$ of $B_n$ sending 
    $\alpha\in A$ to 
    $\kappa_{\gamma}(\alpha):=\gamma \circ \alpha \circ \gamma^{-1}$.
    Obviously,
    we have $\gamma(\alpha g) = \kappa_{\gamma}(\alpha)(\gamma g)$.
    Thus the pair
    $(\kappa_{\gamma}, \gamma)$ is an isomorphism
    from the $A$-set $G$ onto the 
    $(S_n\wr C_2)$-set $S_n$.    
    In particular, $\kappa_{\gamma}$ sends subgroups
    of $A$ which act regularly on $G$, onto subgroups
    of $S_n \wr C_2$ which act regularly on $S_n$.

    The left and right multiplications with elements of $G$
    induce regular subgroups of $A$.
    These are sent to regular subgroups
    $L$ and $R$ (say) of $S_n \wr C_2$.
    Since left and right multiplications centralize each other,
    the subgroups $L$ and $R$ centralize each other.
    If $n\geq 4$, then Lemma~\ref{l:snwrc2_reg} yields
    that $L= S_n \times 1$ or $L=1\times S_n$. 
    Since $L\iso G$, we have that $G\iso S_n$.
    In view of Lemma~\ref{l:snbirk_unique},
    this finishes the proof in case $n\geq 4$.
    
    In the case $n=3$, however, there is one additional
    possibility (up to conjugacy in $S_3 \wr C_2$),
    namely that 
    $L= R = C_2 \times C_3\iso C_6$.
    And indeed, the action of
    $C_2\times C_3$ on $\mat_3(\reals)$ yields the 
    Birkhoff polytope~$B_3$ as orbit polytope of $C_6$,    
    and this orbit polytope is affinely equivalent 
    to the representation polytope $P(D)$,
    where $D\colon C_6 \to \GL(4,\reals)$
    sends a generator of $C_6$ to
    \[ \begin{pmatrix*}[r]
         0 & 1 & & \\
        -1 & -1 & & \\
           &    & 0 & -1 \\
           &    & 1 & 1 
       \end{pmatrix*}.
    \]
\end{proof}  
